\documentclass[11pt, a4paper]{article}
\usepackage[T1]{fontenc}
\usepackage[utf8]{inputenc}
\usepackage[a4paper, top=2.5cm, bottom=2.5cm, left=2.5cm, right=2.5cm]{geometry}
\usepackage{booktabs}
\usepackage{tabularx}
\usepackage{enumitem}
\usepackage[hidelinks]{hyperref}
\usepackage{parskip}
\usepackage{graphicx}
\usepackage{amsmath}
\usepackage{float}
\setlength{\parindent}{0pt}
\pagestyle{plain}

\begin{document}

% ------------------------------------------------------------------------------
% PORTADA
% ------------------------------------------------------------------------------
\begin{center}
{\Large \textbf{Pontificia Universidad Javeriana}}\\[0.3cm]
Facultad de Ciencias Economicas y Administrativas\\[0.3cm]
Estadistica -- 2026-II\\[1cm]
{\LARGE \textbf{Centro vs periferia: ?`la universidad cierra la brecha territorial?}}\\[0.5cm]
{\large Informe Final del Proyecto}\\[1cm]
\textbf{Grupo 7}\\[0.3cm]
Maez Cortes Kevin Santiago\\
Luzz Viloria Damian Mauricio\\
Silva Rojas Julian Felipe\\[1cm]
27 de noviembre de 2026
\end{center}

\newpage
\tableofcontents
\newpage

% ------------------------------------------------------------------------------
\section{Introduccion}
% ------------------------------------------------------------------------------
% Contexto: desigualdad centro-periferia en Colombia.
% Pregunta: los estudiantes de departamentos perifericos (Choco, Guainia,
% Vaupes, Amazonas, Vichada, Guaviare) cierran la brecha con el centro
% (Bogota, Antioquia, Valle) al pasar por la universidad?
% Concepto de valor agregado (delta) como medida de progreso educativo.

% ------------------------------------------------------------------------------
\section{Datos y Metodologia}
% ------------------------------------------------------------------------------
% Fuente: cruce Saber 11 -- Saber Pro (ICFES).
% Archivos: cruce\_saber11\_saberpro.parquet, saber11.parquet, saber\_pro.parquet.
% Variable dependiente: delta = scale(PUNT\_GLOBAL\_sbpro) - scale(PUNT\_GLOBAL\_sb11).
% Construccion de la variable region: Periferia vs Centro
%   a partir de ESTU\_DEPTO\_RESIDE\_sb11.
% Predictores: region, PUNT\_GLOBAL\_sb11, FAMI\_ESTRATOVIVIENDA\_sb11,
%   INST\_ORIGEN\_sbpro.
% Modelo OLS unico con diagnosticos completos.

% ------------------------------------------------------------------------------
\section{Analisis Exploratorio (Fase 1)}
% ------------------------------------------------------------------------------
% Distribucion de delta por region.
% Boxplot del valor agregado centro vs periferia.
% Estadisticas descriptivas por region, estrato e inst\_origen.

% ------------------------------------------------------------------------------
\section{Intervalos de Confianza y Pruebas de Hipotesis (Fase 2)}
% ------------------------------------------------------------------------------
% IC para la media de delta por region.
% Prueba t para diferencia de medias delta centro vs periferia.
% Prueba de Wilcoxon como alternativa no parametrica.

% ------------------------------------------------------------------------------
\section{Modelo de Regresion (Fase 3)}
% ------------------------------------------------------------------------------
% Modelo: delta $\sim$ region + PUNT\_GLOBAL\_sb11 + FAMI\_ESTRATOVIVIENDA\_sb11
%   + INST\_ORIGEN\_sbpro.
% Tabla de coeficientes con IC al 95\%.
% Predicciones: delta predicho para periferia vs centro al P50 de SB11.
% El modelo predice que los estudiantes de periferia cierran la brecha?

% ------------------------------------------------------------------------------
\section{Diagnosticos y Robustez}
% ------------------------------------------------------------------------------
% Graficos de residuales (residuales vs ajustados, QQ-plot, Scale-Location).
% Test de Breusch-Pagan (homocedasticidad).
% VIF (multicolinealidad).
% Distancia de Cook (observaciones influyentes).
% Errores robustos HC1: comparacion con errores estandar clasicos.

% ------------------------------------------------------------------------------
\section{Resultados y Discusion}
% ------------------------------------------------------------------------------
% Magnitud y significancia del coeficiente de region.
% Comparacion de predicciones centro vs periferia.
% Scatter SB11 vs delta por region.
% Reflexion sobre el papel de la universidad en la equidad territorial.

% ------------------------------------------------------------------------------
\section{Conclusiones}
% ------------------------------------------------------------------------------
% Respuesta a la pregunta de investigacion.
% Limitaciones del estudio (tamano muestral de departamentos perifericos).
% Implicaciones para politica de descentralizacion educativa.

% ------------------------------------------------------------------------------
\section*{Referencias}
% ------------------------------------------------------------------------------
% ICFES, literatura sobre desigualdad territorial y educacion en Colombia.

% ------------------------------------------------------------------------------
\section*{Anexo: Codigo R}
% ------------------------------------------------------------------------------
% Incluir codigo R utilizado (fase3\_grupo7.R) o referenciarlo.

\end{document}
